\begin{table}[htbp]
\floatleft
\caption{The Three Ontological State Spaces and Their Production Frames}
\label{tab:spaces}
\vspace{2pt}
\begingroup\small
\begin{tabularx}{\linewidth}{@{}p{0.13\linewidth}YYp{0.17\linewidth}@{}}
\toprule
State space & Hierarchical structure and metadata & Production-frame cardinality & Role in a
scenario \\
\midrule
Profile & Recursive tree of individual-difference constructs: demographics, life context,
five-factor facets, political psychology, ideology, moral foundations, socioeconomics,
information behaviour & 272 categorical variables / 1,467 options; 41 continuous groups / 181
scale leaves; 26 ordinal; 6 trait scalars; 7 identifiers ($\approx$159 features analysed) & The
simulated person whose attributes may moderate susceptibility \\[2pt]
Opinion & Issue Position Taxonomy of directional policy positions nested in issue-domain
clusters; each leaf carries $d_\ell \in \{-1,+1\}$ & 106 directional leaves in 7 clusters (6
analysed); neutral leaves excluded from scoring & The target moved by the attack and scored for
direction-aware effectivity \\[2pt]
Attack & DISARM-red Plan, Prepare, and Execute phases, each a catalogue of tactics and
techniques, with a complexity tier, signal score, and inclusion route & 48,991 filtered triplets
(17,941 distinct leaves) from a $2.2\times10^{11}$ Cartesian upper bound & The adversarial
operation specification presented to the agent \\
\bottomrule
\end{tabularx}
\endgroup
\apanote{Each ontology is a recursively populated tree with stable paths and leaf-level
metadata. The full attack ontology is dual-use and is not reproduced. Cardinalities are the
production frames before maximum-entropy sampling into scenarios.}
\end{table}
